\begin{figure}[tbp]
\centering
\begin{minipage}[t]{.42\linewidth}
\centering\ViewHeading{(a) Reverse divisibility}
\begin{tikzpicture}
\node[vbox] (six) at (0,0) {$6$};
\node[vbox] (two) at (-1.7,-1.5) {$2$};
\node[vbox] (three) at (1.7,-1.5) {$3$};
\node[vresult] (one) at (0,-3) {$1$};
\draw[varrow] (six)--(two) node[midway,above left,font=\footnotesize] {$T_{2,3}$};
\draw[varrow] (six)--(three) node[midway,above right,font=\footnotesize] {$T_{3,2}$};
\draw[varrow] (two)--(one) node[midway,below left,font=\footnotesize] {$T_{1,2}$};
\draw[varrow] (three)--(one) node[midway,below right,font=\footnotesize] {$T_{1,3}$};
\node[font=\footnotesize] at (0,-3.85) {both composites are $T_{1,6}$};
\end{tikzpicture}
\end{minipage}\hfill
\begin{minipage}[t]{.55\linewidth}
\centering\ViewHeading{(b) Boundary and norm compatibility}
\begin{tikzpicture}
\node (sourceone) at (0,0) {$\mathcal C_x(kq)_1$};
\node (sourcezero) at (4.3,0) {$\C$};
\node (targetone) at (0,-1.65) {$\mathcal C_x(q)_1$};
\node (targetzero) at (4.3,-1.65) {$\C$};
\draw[varrow] (sourceone)--(sourcezero) node[midway,above] {$\partial_{kq,x}$};
\draw[varrow] (targetone)--(targetzero) node[midway,below] {$\partial_{q,x}$};
\draw[varrow] (sourceone)--(targetone) node[midway,left] {$T_{q,k}$};
\draw[varrow] (sourcezero)--(targetzero) node[midway,right] {$\id$};
\node[vinput,text width=6.2cm] at (2.05,-3.25)
 {$\|T_{q,k}\|_{2\to2}\leq k$\\$\|T_{q,k}c\|_q\leq qk\|c\|_2=\|c\|_{kq}$};
\end{tikzpicture}
\end{minipage}
\caption{Arithmetic maps point from a multiple to a divisor; they are not
inclusions of a filtration. Weighting the degree-one norm at each node
by its scale makes the maps nonexpansive. The degree-zero map fixes the
marked chain $1$.}
\label{fig:arithmetic-diagram}
\end{figure}